%%%%%%%%%%%%%%%%%%%%%%%%%%%%%%%%%%%%%%%%%%%%%%%%%%%%%%%%%%%%%%%%%%%%%%%%
% APPENDIX
%%%%%%%%%%%%%%%%%%%%%%%%%%%%%%%%%%%%%%%%%%%%%%%%%%%%%%%%%%%%%%%%%%%%%%%%

%%%%%%%%%%%%%%%%%%%%%%%%%%%%%%%%%%%%%%%%%%%%%%%%%%%%%%%%%%%%%%%%%%%%%%%%
\chapter{Use of Generative Artificial Intelligence}
\label{appendix:ai_usage}
%%%%%%%%%%%%%%%%%%%%%%%%%%%%%%%%%%%%%%%%%%%%%%%%%%%%%%%%%%%%%%%%%%%%%%%%

This appendix documents the technical configuration used to integrate generative \gls{ai} assistance into the thesis writing workflow.
The complete configuration is available in the thesis repository at:

\begin{center}
\url{https://github.com/simonmeoni/thesis}
\end{center}

\section{Claude Code Configuration}

This thesis was written with assistance from Claude Code, Anthropic's command-line interface for Claude.
The tool was configured through a \texttt{.claude/} directory at the root of the repository, containing project-specific instructions and custom skills.

\subsection{Project Instructions (CLAUDE.md)}

A \texttt{CLAUDE.md} file at the repository root provides Claude with context about the thesis structure, LaTeX conventions, and writing guidelines.
This file includes:
\begin{itemize}
    \item Overview of the thesis structure and chapter organization
    \item LaTeX conventions (citation styles, acronym usage, figure references)
    \item Writing rules specific to academic thesis formatting
    \item Custom terminology and key concepts
\end{itemize}

\subsection{Custom Skills}

Custom skills were developed to automate repetitive tasks in the thesis workflow.
Each skill is defined in a markdown file within \texttt{.claude/skills/} and provides specialized instructions for specific tasks:

\begin{description}
    \item[\texttt{/compile}] Compiles the thesis using \texttt{latexmk} and reports errors and warnings in a structured format.
    \item[\texttt{/word-count}] Counts words in thesis chapters using \texttt{texcount} and provides statistics by chapter.
    \item[\texttt{/bib-check}] Checks \texttt{bibliography.bib} for syntax errors, duplicate entries, and undefined citations.
    \item[\texttt{/cite-assist}] Verifies citations during writing, checks bibliography locally, searches online for missing references, and fixes potential hallucinations.
    \item[\texttt{/check-todos}] Finds TODO and FIXME comments in LaTeX files to track remaining tasks.
    \item[\texttt{/review}] Performs comprehensive review of chapter structure, content quality, and LaTeX best practices.
    \item[\texttt{/biblio-review}] Critical review of bibliography content, coverage, and relevance for a given chapter or topic.
    \item[\texttt{/web-search}] Searches the web for research papers and documentation to support literature review.
    \item[\texttt{/commit}] Creates conventional commits with proper format following the \texttt{type(scope): description} convention.
    \item[\texttt{/push}] Pushes commits to the remote repository.
\end{description}

\subsection{Workflow Integration}

The AI-assisted workflow followed these principles:
\begin{enumerate}
    \item \textbf{Human-driven content}: All scientific ideas, arguments, and conclusions originated from the author.
    \item \textbf{AI-assisted refinement}: Claude was used to improve clarity, fix LaTeX errors, and suggest structural improvements.
    \item \textbf{Critical review}: All \gls{ai} suggestions were reviewed and edited before inclusion.
    \item \textbf{Version control}: All changes were tracked through git commits, enabling full traceability of modifications.
\end{enumerate}

\subsection{Transparency}

The \texttt{.claude/} directory is included in the public repository, allowing readers to inspect the exact configuration used during thesis preparation.
Commit messages include co-authorship attribution (\texttt{Co-Authored-By: Claude}) when \gls{ai} assistance was used for the changes.

%%%%%%%%%%%%%%%%%%%%%%%%%%%%%%%%%%%%%%%%%%%%%%%%%%%%%%%%%%%%%%%%%%%%%%%%
\chapter{Weak Supervision: Detailed Results}
\label{appendix:weak_supervision}
%%%%%%%%%%%%%%%%%%%%%%%%%%%%%%%%%%%%%%%%%%%%%%%%%%%%%%%%%%%%%%%%%%%%%%%%

This appendix provides detailed experimental results for the weak supervision and knowledge distillation study presented in Chapter 4.

\section{Detailed Results by Language and Model}

\subsection{Monolingual Silver Annotations (r = 0 or 1)}
\label{appendix:setting_a}

Table \ref{tab:appendix_setting_a} presents F1-scores for \Ssilver with $\rmix \in \{0, 1\}$ across all evaluated languages and models, comparing \distilmodels ($\rmix=1$) and \weakModel ($\rmix=0$) approaches.

\begin{table}[h]
    \centering
    \small
    \resizebox{1\textwidth}{!}{
    \begin{tabular}{@{}cccc@{}}%
\toprule%
Language&Model&\multicolumn{2}{c}{F1-Score}\\%
&&$r=1$&$r=0$\\%
\midrule%
en&\verb|emilyalsentzer/Bio_ClinicalBERT|&0.66  $\pm$ 0.01&\textbf{0.68} \bm{$\pm$} \textbf{0.01}\\%
&\verb|roberta-base|&\textbf{0.67} \bm{$\pm$} \textbf{0.01}&0.65  $\pm$ 0.01\\%
&\verb|xlm-roberta-base|&\textbf{0.66} \bm{$\pm$} \textbf{0.01}&0.65  $\pm$ 0.01\\%
\midrule%
es&\verb|BSC-LT/roberta-base-biomedical-es|&0.72  $\pm$ 0.01&\textbf{0.78} \bm{$\pm$} \textbf{0.01}\\%
&\verb|dccuchile/bert-base-spanish-wwm-cased|&0.69  $\pm$ 0.01&\textbf{0.76} \bm{$\pm$} \textbf{0.01}\\%
&\verb|xlm-roberta-base|&0.68  $\pm$ 0.02&\textbf{0.73} \bm{$\pm$} \textbf{0.01}\\%
\midrule%
eu&\verb|ixa-ehu/berteus-base-cased|&\textbf{0.60} \bm{$\pm$} \textbf{0.03}&0.54  $\pm$ 0.03\\%
&\verb|xlm-roberta-base|&\textbf{0.61} \bm{$\pm$} \textbf{0.04}&0.47  $\pm$ 0.01\\%
\midrule%
fr&\verb|Dr-BERT/DrBERT-7GB|&\textbf{0.74} \bm{$\pm$} \textbf{0.01}&0.70  $\pm$ 0.01\\%
&\verb|camembert-base|&\textbf{0.75} \bm{$\pm$} \textbf{0.01}&0.69  $\pm$ 0.04\\%
&\verb|xlm-roberta-base|&\textbf{0.74} \bm{$\pm$} \textbf{0.01}&0.72  $\pm$ 0.02\\%
\midrule%
it&\verb|dbmdz/bert-base-italian-cased|&\textbf{0.74} \bm{$\pm$} \textbf{0.01}&0.73  $\pm$ 0.00\\%
&\verb|xlm-roberta-base|&\textbf{0.75} \bm{$\pm$} \textbf{0.01}&0.72  $\pm$ 0.02\\\bottomrule%
%
\end{tabular}
    }
    \caption{F1-scores for \Ssilver with $\rmix \in \{0, 1\}$ by language and model. Comparison between \distilmodels ($\rmix=1$) and \weakModel ($\rmix=0$).}
    \label{tab:appendix_setting_a}
\end{table}

\clearpage

\subsection{Monolingual Silver Validation}
\label{appendix:setting_b}

Table \ref{tab:appendix_setting_b} shows results for \Sval, where models are trained on the manually corrected subset.

\begin{table}[h]
    \centering
    \small
    \resizebox{1\textwidth}{!}{
        \begin{tabular}{@{}cccc@{}}%
\toprule%
Language&Model&\multicolumn{2}{c}{F1-Score}\\%
&&$\rmix=1$&$\rmix=0$\\%
\midrule%
en&\verb|emilyalsentzer/Bio_ClinicalBERT|&0.60&\textbf{0.65}\\%
&\verb|roberta-base|&0.61&\textbf{0.70}\\%
&\verb|xlm-roberta-base|&0.43&\textbf{0.69}\\%
\midrule%
es&\verb|BSC-LT/roberta-base-biomedical-es|&0.71&\textbf{0.78}\\%
&\verb|dccuchile/bert-base-spanish-wwm-cased|&0.70&\textbf{0.77}\\%
&\verb|xlm-roberta-base|&0.62&\textbf{0.73}\\%
\midrule%
eu&\verb|ixa-ehu/berteus-base-cased|&0.54&\textbf{0.72}\\%
&\verb|xlm-roberta-base|&0.51&\textbf{0.68}\\%
\midrule%
fr&\verb|Dr-BERT/DrBERT-7GB|&0.73&\textbf{0.75}\\%
&\verb|camembert-base|&\textbf{0.71}&0.60\\%
&\verb|xlm-roberta-base|&0.68&\textbf{0.68}\\%
\midrule%
it&\verb|dbmdz/bert-base-italian-cased|&0.63&\textbf{0.70}\\%
&\verb|xlm-roberta-base|&0.55&\textbf{0.57}\\\bottomrule%
%
\end{tabular}
    }
    \caption{F1-scores for \Sval by language and model.}
    \label{tab:appendix_setting_b}
\end{table}

\clearpage

\subsection{Monolingual Mixed-Supervised}
\label{appendix:setting_c}

Table \ref{tab:appendix_setting_c} presents results for \Sstar (i.e., the combined \Ssilver + \Sval training set), using \hybridmodel approach that combines \Ssilver annotations with manually corrected \Sval data.

\begin{table}[h]
    \centering
    \small
    \resizebox{1\textwidth}{!}{
        \begin{tabular}{@{}cccc@{}}%
\toprule%
Language&Model&\multicolumn{2}{c}{F1-Score}\\%
&&$\rmix=1$&$\rmix=0$\\%
\midrule%
en&\verb|emilyalsentzer/Bio_ClinicalBERT|&\textbf{0.70} \bm{$\pm$} \textbf{0.01}&0.69  $\pm$ 0.01\\%
&\verb|roberta-base|&\textbf{0.70} \bm{$\pm$} \textbf{0.01}&0.68  $\pm$ 0.01\\%
&\verb|xlm-roberta-base|&\textbf{0.67} \bm{$\pm$} \textbf{0.01}&0.67  $\pm$ 0.01\\%
\midrule%
es&\verb|BSC-LT/roberta-base-biomedical-es|&0.77  $\pm$ 0.02&\textbf{0.80} \bm{$\pm$} \textbf{0.01}\\%
&\verb|dccuchile/bert-base-spanish-wwm-cased|&0.76  $\pm$ 0.01&\textbf{0.78} \bm{$\pm$} \textbf{0.00}\\%
&\verb|xlm-roberta-base|&0.76  $\pm$ 0.01&\textbf{0.76} \bm{$\pm$} \textbf{0.01}\\%
\midrule%
eu&\verb|ixa-ehu/berteus-base-cased|&0.70  $\pm$ 0.03&\textbf{0.72} \bm{$\pm$} \textbf{0.02}\\%
&\verb|xlm-roberta-base|&\textbf{0.70} \bm{$\pm$} \textbf{0.01}&0.57  $\pm$ 0.09\\%
\midrule%
fr&\verb|Dr-BERT/DrBERT-7GB|&\textbf{0.75} \bm{$\pm$} \textbf{0.01}&0.73  $\pm$ 0.02\\%
&\verb|camembert-base|&\textbf{0.76} \bm{$\pm$} \textbf{0.01}&0.74  $\pm$ 0.01\\%
&\verb|xlm-roberta-base|&\textbf{0.74} \bm{$\pm$} \textbf{0.01}&0.69  $\pm$ 0.04\\%
\midrule%
it&\verb|dbmdz/bert-base-italian-cased|&\textbf{0.75} \bm{$\pm$} \textbf{0.01}&0.75  $\pm$ 0.01\\%
&\verb|xlm-roberta-base|&\textbf{0.74} \bm{$\pm$} \textbf{0.00}&0.74  $\pm$ 0.01\\\bottomrule%
%
\end{tabular}
    }
    \caption{F1-scores for \Sstar by language and model, comparing \distilmodels ($\rmix=1$) and \weakModel ($\rmix=0$).}
    \label{tab:appendix_setting_c}
\end{table}

\clearpage

\subsection{Monolingual Silver: Detailed Results}
\label{appendix:setting_d}

Table \ref{tab:appendix_setting_d} provides comprehensive results for \Ssilver (\mono), including F1-score, Precision, and Recall for different mixing ratios $\rmix \in \{0, r_{max}, 1\}$.

\begin{landscape}
\begin{table}[p]
    \centering
    \small
    \resizebox{1.5\textwidth}{!}{
    \begin{tabular}{@{}cccccccccccc@{}}%
\toprule%
Language&Model&$r_{max}$&\multicolumn{3}{c}{$r=1$}&\multicolumn{3}{c}{$r_{max}$}&\multicolumn{3}{c}{$r=0$}\\%
&&&F1-Score&Precision&Recall&F1-Score&Precision&Recall&F1-Score&Precision&Recall\\%
\midrule%
en&\verb|xlm-roberta-base|&0.5&0.65 \bm{$\pm$} 0.01&0.72 \bm{$\pm$} 0.01&0.63 \bm{$\pm$} 0.01&0.71 \bm{$\pm$} 0.01&0.71 \bm{$\pm$} 0.03&0.71 \bm{$\pm$} 0.02&0.66 \bm{$\pm$} 0.01&0.68 \bm{$\pm$} 0.03&0.64 \bm{$\pm$} 0.02\\%
&\verb|roberta-base|&0.4&0.65 \bm{$\pm$} 0.01&0.72 \bm{$\pm$} 0.01&0.62 \bm{$\pm$} 0.02&0.71 \bm{$\pm$} 0.01&0.72 \bm{$\pm$} 0.01&0.71 \bm{$\pm$} 0.03&0.67 \bm{$\pm$} 0.01&0.70 \bm{$\pm$} 0.01&0.65 \bm{$\pm$} 0.01\\%
&\verb|emilyalsentzer/Bio_ClinicalBERT|&0.6&0.68 \bm{$\pm$} 0.01&0.73 \bm{$\pm$} 0.01&0.66 \bm{$\pm$} 0.02&0.72 \bm{$\pm$} 0.01&0.75 \bm{$\pm$} 0.02&0.70 \bm{$\pm$} 0.02&0.66 \bm{$\pm$} 0.01&0.71 \bm{$\pm$} 0.02&0.63 \bm{$\pm$} 0.01\\%
\midrule%
es&\verb|xlm-roberta-base|&0.4&0.73 \bm{$\pm$} 0.01&0.80 \bm{$\pm$} 0.01&0.70 \bm{$\pm$} 0.01&0.76 \bm{$\pm$} 0.01&0.77 \bm{$\pm$} 0.02&0.76 \bm{$\pm$} 0.02&0.68 \bm{$\pm$} 0.02&0.74 \bm{$\pm$} 0.02&0.65 \bm{$\pm$} 0.04\\%
&\verb|dccuchile/bert-base-spanish-wwm-cased|&0.4&0.76 \bm{$\pm$} 0.01&0.81 \bm{$\pm$} 0.01&0.73 \bm{$\pm$} 0.02&0.78 \bm{$\pm$} 0.00&0.80 \bm{$\pm$} 0.02&0.76 \bm{$\pm$} 0.02&0.69 \bm{$\pm$} 0.01&0.72 \bm{$\pm$} 0.02&0.66 \bm{$\pm$} 0.03\\%
&\verb|BSC-LT/roberta-base-biomedical-es|&0.4&0.78 \bm{$\pm$} 0.01&0.82 \bm{$\pm$} 0.01&0.75 \bm{$\pm$} 0.01&0.79 \bm{$\pm$} 0.01&0.80 \bm{$\pm$} 0.04&0.79 \bm{$\pm$} 0.03&0.72 \bm{$\pm$} 0.01&0.75 \bm{$\pm$} 0.02&0.69 \bm{$\pm$} 0.02\\%
\midrule%
eu&\verb|xlm-roberta-base|&1&0.47 \bm{$\pm$} 0.01&0.63 \bm{$\pm$} 0.14&0.44 \bm{$\pm$} 0.01&{-}&{-}&{-}&0.61 \bm{$\pm$} 0.04&0.73 \bm{$\pm$} 0.02&0.56 \bm{$\pm$} 0.05\\%
&\verb|ixa-ehu/berteus-base-cased|&1&0.54 \bm{$\pm$} 0.03&0.91 \bm{$\pm$} 0.01&0.49 \bm{$\pm$} 0.02&{-}&{-}&{-}&0.60 \bm{$\pm$} 0.03&0.74 \bm{$\pm$} 0.02&0.54 \bm{$\pm$} 0.03\\%
\midrule%
fr&\verb|xlm-roberta-base|&0.4&0.72 \bm{$\pm$} 0.02&0.80 \bm{$\pm$} 0.02&0.67 \bm{$\pm$} 0.03&0.75 \bm{$\pm$} 0.01&0.75 \bm{$\pm$} 0.02&0.75 \bm{$\pm$} 0.02&0.74 \bm{$\pm$} 0.01&0.74 \bm{$\pm$} 0.02&0.74 \bm{$\pm$} 0.02\\%
&\verb|camembert-base|&0.5&0.74 \bm{$\pm$} 0.01&0.69 \bm{$\pm$} 0.01&0.83 \bm{$\pm$} 0.01&0.76 \bm{$\pm$} 0.01&0.73 \bm{$\pm$} 0.02&0.81 \bm{$\pm$} 0.03&0.75 \bm{$\pm$} 0.00&0.82 \bm{$\pm$} 0.01&0.71 \bm{$\pm$} 0.01\\%
&\verb|Dr-BERT/DrBERT-7GB|&0.6&0.70 \bm{$\pm$} 0.01&0.84 \bm{$\pm$} 0.00&0.64 \bm{$\pm$} 0.01&0.76 \bm{$\pm$} 0.01&0.76 \bm{$\pm$} 0.01&0.77 \bm{$\pm$} 0.01&0.74 \bm{$\pm$} 0.01&0.77 \bm{$\pm$} 0.04&0.73 \bm{$\pm$} 0.04\\%
\midrule%
it&\verb|xlm-roberta-base|&0.6&0.72 \bm{$\pm$} 0.02&0.78 \bm{$\pm$} 0.02&0.70 \bm{$\pm$} 0.04&0.75 \bm{$\pm$} 0.02&0.76 \bm{$\pm$} 0.02&0.76 \bm{$\pm$} 0.02&0.75 \bm{$\pm$} 0.01&0.75 \bm{$\pm$} 0.02&0.75 \bm{$\pm$} 0.02\\%
&\verb|dbmdz/bert-base-italian-cased|&0.8&0.73 \bm{$\pm$} 0.00&0.76 \bm{$\pm$} 0.02&0.73 \bm{$\pm$} 0.02&0.75 \bm{$\pm$} 0.00&0.72 \bm{$\pm$} 0.02&0.81 \bm{$\pm$} 0.02&0.74 \bm{$\pm$} 0.01&0.70 \bm{$\pm$} 0.02&0.81 \bm{$\pm$} 0.02\\\bottomrule%
%
\end{tabular}
    }
    \caption{Detailed results for \Ssilver by language and model. Shows performance at $\rmix=1$ (\distilmodels only), $\rmix=0$ (\weakModel only), and optimal ratio $\rmix_{max}$ (\hybridmodel).}
    \label{tab:appendix_setting_d}
\end{table}
\end{landscape}
\clearpage

\subsection{Multilingual Silver: Detailed Results}
\label{appendix:setting_e}

Table \ref{tab:appendix_setting_e} provides comprehensive results for \Smulti using multilingual training with \verb|xlm-roberta-base|.

\begin{landscape}
\begin{table}[p]
    \centering
    \small
    \resizebox{1.5\textwidth}{!}{
    \begin{tabular}{@{}cccccc@{}}%
\toprule%
Language&$\rmix_{max}$&Metric&{$\rmix=1$}&{$\rmix_{max}$}&{$\rmix=0$}\\%
\midrule%
\multirow{3}{*}{en}&\multirow{3}{*}{1}&F1-Score&0.64 \bm{$\pm$} 0.01&{-}&0.66 \bm{$\pm$} 0.03\\%
&&Precision&0.74 \bm{$\pm$} 0.01&{-}&0.75 \bm{$\pm$} 0.03\\%
&&Recall&0.60 \bm{$\pm$} 0.01&{-}&0.62 \bm{$\pm$} 0.05\\%
\midrule%
\multirow{3}{*}{es}&\multirow{3}{*}{0.6}&F1-Score&0.71 \bm{$\pm$} 0.01&0.74 \bm{$\pm$} 0.01&0.71 \bm{$\pm$} 0.01\\%
&&Precision&0.82 \bm{$\pm$} 0.01&0.83 \bm{$\pm$} 0.02&0.80 \bm{$\pm$} 0.03\\%
&&Recall&0.65 \bm{$\pm$} 0.02&0.69 \bm{$\pm$} 0.02&0.65 \bm{$\pm$} 0.02\\%
\midrule%
\multirow{3}{*}{eu}&\multirow{3}{*}{0.5}&F1-Score&0.57 \bm{$\pm$} 0.02&0.59 \bm{$\pm$} 0.02&0.57 \bm{$\pm$} 0.01\\%
&&Precision&0.86 \bm{$\pm$} 0.02&0.79 \bm{$\pm$} 0.02&0.76 \bm{$\pm$} 0.04\\%
&&Recall&0.51 \bm{$\pm$} 0.02&0.54 \bm{$\pm$} 0.02&0.51 \bm{$\pm$} 0.01\\%
\midrule%
\multirow{3}{*}{fr}&\multirow{3}{*}{0.6}&F1-Score&0.69 \bm{$\pm$} 0.02&0.71 \bm{$\pm$} 0.01&0.70 \bm{$\pm$} 0.02\\%
&&Precision&0.83 \bm{$\pm$} 0.01&0.81 \bm{$\pm$} 0.02&0.80 \bm{$\pm$} 0.02\\%
&&Recall&0.62 \bm{$\pm$} 0.02&0.65 \bm{$\pm$} 0.02&0.64 \bm{$\pm$} 0.03\\%
\midrule%
\multirow{3}{*}{it}&\multirow{3}{*}{0.8}&F1-Score&0.73 \bm{$\pm$} 0.01&0.76 \bm{$\pm$} 0.01&0.76 \bm{$\pm$} 0.01\\%
&&Precision&0.78 \bm{$\pm$} 0.01&0.76 \bm{$\pm$} 0.01&0.75 \bm{$\pm$} 0.03\\%
&&Recall&0.71 \bm{$\pm$} 0.01&0.77 \bm{$\pm$} 0.02&0.78 \bm{$\pm$} 0.02\\\bottomrule%
%
\end{tabular}

    }
    \caption{Detailed results for \Smulti (multilingual training). Shows performance at $\rmix=1$ (\distilmodels only), $\rmix=0$ (\weakModel only), and optimal ratio $\rmix_{max}$ (\hybridmodel).}
    \label{tab:appendix_setting_e}
\end{table}
\end{landscape}

\clearpage


